% Original PGFPlots regression case for the PGF constant and jump handlers.
% The 6cm x 4cm box has one cm per data unit. Neither trace samples x=0 or 6.
% Blue holds y=2 at the left edge and y=2.5 at the right edge. Its y=5
% plateau is invisible; only the risers at x=3 and 4 reach the top edge.
% Red holds y=.5,1.5,.75 over visible x intervals [0,.5], [.5,3.5],
% [4.5,6]. The y=-1 interval [3.5,4.5] leaves a gap, not a bottom-edge line.
\documentclass[border=2pt]{standalone}
\usepackage{pgfplots}
\pgfplotsset{compat=1.18}
\begin{document}
\begin{tikzpicture}
  \begin{axis}[width=6cm,height=4cm,scale only axis,axis lines=box,
    xmin=0,xmax=6,ymin=0,ymax=4,enlargelimits=false,clip=true,ticks=none]
    \addplot[blue,thick,const plot mark left,mark=none]
      coordinates {(-2,2) (1,3) (3,5) (4,2.5) (8,1.5)};
    \addplot[red,thick,jump mark mid,mark=none]
      coordinates {(-2,0.5) (3,1.5) (4,-1) (5,0.75) (8,1.25)};
  \end{axis}
\end{tikzpicture}
\end{document}
